\scanpage{25}
% archive sheet number: 24 (inverted)
\mspage{-I- 1 -}\hfill \handwritten{n°383}

\S~1. \uuline{Résultats préliminaires}.

1. \uline{Lemmes sur les chambres}.

Soit $V$ un espace vectoriel réel de dimension finie.

DEFINITION 1.- \uline{Une partie} $C$ \uline{de} $V$ \uline{est appelée une chambre si elle est non-vide et s'il existe une famille localement finie} $\underline H$ (Top. Gén., I, \handwritten{3}ème éd., \S~1, n°\handwritten{5}, déf.8) \uline{d'hyperplans affines de} $V$ \uline{et}, \uline{pour tout} $H\in\underline H$, \uline{un demi-espace ouvert} $D_H(C)$ \uline{défini par} $H$ \uline{tels que} $C=\displaystyle\bigcap_{H\in\underline H}D_H(C)$.

Un point quelconque $x$ de $C$ n'est dans aucun hyperplan $H\in\underline H$ ; comme $\underline H$ est localement finie, il existe une boule de centre $x$ qui ne rencontre aucun $H\in\underline H$, et qui est donc contenue dans chaque demi\unsure{-|.}espace $D_H(C)$ ; il en résulte que toute chambre est \uline{ouverte}. Il est clair que toute chambre est \uline{convexe}.

DEFINITION 2.- \uline{Soit} $C$ \uline{une chambre}. \uline{Un hyperplan affine} $L$ \uline{de} $V$ \uline{est appelé un mur de} $C$ \uline{s'il existe un demi-espace ouvert} $L^+$ \uline{défini par} $L$, \uline{un point} $x$ \uline{de} $L$ \uline{et un voisinage} $U$ \uline{de} $x$ \uline{tels que} $U\cap L^+=U\cap C$.

LEMME 1.- \uline{Soient} $L,L'$ \uline{deux hyperplans affines de} $V$, $L^+$ \uline{et} $L^-$ (resp.\ $L^{\prime+}$ \uline{et} $L^{\prime-}$) \uline{les demi-espaces ouverts définis par} $L$ (resp.\ $L'$). \uline{S'il existe un point} $x$ \uline{de} $L\cap L'$ \uline{et un voisinage} $U$ \uline{de} $x$ \uline{tels que} $U\cap L^+\subset L^{\prime+}$, \uline{alors} $L^+=L^{\prime+}$, $L^-=L^{\prime-}$ \uline{et} $L=L'$.

Par translation on se ramène au cas où $x=0$. Comme tout élément de $L^+$ se déduit \margin{\handwritten{\unsure{$>$}}} d'un élément de $U\cap L^+$ par une homothétie de rapport $>0$, on a $L^+\subset L^{\prime+}$ ; d'où, par symétrie, $L^-\subset L^{\prime-}$, et donc $L^+\cup L^-\subset L^{\prime+}\cap L^{\prime-}$. Par passage aux complémentaires on en déduit $L'\subset L$, d'où $L'=L$ car ce sont des hyperplans. Par conséquent $L^+=L^{\prime+}$ et $L^-=L^{\prime-}$.

LEMME 2.- \uline{Soient} $C$ \uline{une chambre de} $V$, $\underline H$ \uline{une famille localement finie d'hyperplans telle que} $C=\displaystyle\bigcap_{H\in\underline H}D_H(C)$, \uline{et} $\underline H'$ \uline{la famille des murs de} $C$. \uline{Alors} :
